\section{Release used, records and the evaluated split}
\label{app:release}

\paragraph{The April 2026 release and Zenodo v11.}
The segmentation experiments ran on the April 2026 release
(\S\ref{sec:dataset}). Its separate spine-only and
pelvic-only records, for patients whose two source
annotations lie on different acquisitions, are what the
partial-annotation protocol (\S\ref{sec:partial}) trains
on. The current archived version, v11, differs from it as
follows: one record per patient (802 records), each
separate-mode patient's missing pelvis completed on the
spine's series by out-of-fold pseudolabelling; VerSe-native
identifiers (L1--L6 = 20--25, sacrum 26, hips 30/31) with
femora, per-level ribs, lumbar ribs and surgical hardware
added as further classes; hip-laterality corrections on 22
records and hand corrections on five; both Castellvi reads
and their consensus in the manifest; and the CT images
supplied by TCIA series crosswalk and rebuild script, with
a HuggingFace mirror carrying the volumes. The 802 patients
and the 33 Castellvi-graded patients are the same in both.

\paragraph{Records and split.}
\label{app:split}
Table~\ref{tab:scan_stats} gives the composition and the
record sets the experiments used. The train, validation
and test manifests are those of the release's 1 May 2026
revision, redrawn two days before training began; the
converter sends the train and validation manifests to the
cross-validation pool and holds out the test manifest.
The five folds are the release's patient-level fold
assignment copied onto the pool's records, so within the
pool no patient spans two validation folds. Each
validation-manifest record therefore trains four of the
five fold models, and each fold's best checkpoint is
selected on that fold's validation records, never on the
test manifest.

\begin{table}[h]
\caption{\dataset{} composition (April 2026 release) and
the record sets of the experiments.}
\label{tab:scan_stats}
\centering
\small
\begin{tabular}{lrr}
\toprule
\bfseries Set & \bfseries Patients & \bfseries Records \\
\midrule
Fused (spine + pelvis)                         & 342 & 342 \\
Separate-mode (spine record + pelvic record)   & 351 & 702 \\
Spine-only / pelvic-only                       & 89 / 20 & 89 / 20 \\
All                                            & 802 & 1{,}153 \\
\quad with spine / pelvic labels               & --- & 782 / 713 \\
\quad of the 33 LSTV patients                  & 33 & 53 \\
\midrule
Train manifest                                 & 639 & 805 \\
Validation manifest                            & 170 & 174 \\
Cross-validation pool (train + validation)     & --- & 979 \\
Test manifest (held out)                       & 167 & 174 \\
\bottomrule
\end{tabular}
\end{table}

The test manifest is a record-level split stratified by
LSTV class and match type. Of its 174 records (167
patients), 82 records from 75 patients are
patient-disjoint from the pool: 52 fused and 16
single-mode patients (14 spine-only, 2 pelvic-only), and
7 separate-mode patients with both records held out. The
other 92 are separate-mode records (52 pelvic, 40 spine)
whose sibling---the patient's other acquisition, labelled
only for the other region---is in the pool (81 from the
train manifest, 11 from the validation manifest) and
trains four of the five fold models. Among the 33 LSTV
patients, 167, 672 and 125 are held out entirely; 401,
555, 123 and 4 have an LSTV-tagged held-out record and a
pool record (tagged normal for 401, 555 and 123 by the
refinement rule of \S\ref{sec:partial}, ambiguous for 4);
104 has its spine record (tagged sacr\_count) in the pool
and its pelvic record, tagged normal, held out. The pool
siblings of 104 and 4 carry an LSTV tag and were
therefore oversampled in the folds that trained on them.
By record, the 33 LSTV patients' 53 records are 44 in the
pool (28 LSTV-tagged) and 9 held out (7 LSTV-tagged).
Table~\ref{tab:castellvi_dist} counts patients: the Pool
and Test columns overlap in one patient (4) by tag, and in
five (401, 555, 123, 4, 104) by any record.

\paragraph{The Normal stratum.}
The Normal rows of Tables~\ref{tab:comparison}
and~\ref{tab:baseline} hold 167 records from 161
patients: 165 records of 159 normal patients and the
pelvic records of LSTV patients 104 (count-style
sacralization) and 672 (lumbarization), tagged normal
because they carry no spine labels. Removing those two
records moves the Normal ensemble sacrum mean from 0.9719
to 0.9717, the merged hip from 0.9884 to 0.9882 and the
mean foreground Dice from 0.9658 to 0.9655; spine classes
and junction-DSC are unaffected.

\paragraph{What is public.}
Neither the April release's 1{,}153 records nor the
evaluated split is at a location this paper cites. The
Zenodo archive (v11) and the OpenSpineConsortium
HuggingFace mirror hold the 802-record v11 layout, and the
archive's \texttt{splits\_5fold.json} is the patient-level
fold assignment over the 802 patients, with no test split.
Its per-patient subtype map is not the six-way taxonomy of
this paper: like the unmerged control's
(Appendix~\ref{app:unmerged}), it has no count-style class
and tags one more patient lumbarization (lumbarization 15,
sacralization 15, semi-sacralization 2, ambiguous 2, normal
768), so strata drawn from it differ from those of
Tables~\ref{tab:sota_stratified}
and~\ref{tab:castellvi_dist}. The six-way subtypes are the
\texttt{lstv\_phenotype} field of the v11 manifest
(lumbarization 14, count-style sacralization 9, full
sacralization 6, semi-sacralization 2, ambiguous 2, empty
for the 769 normal patients).
The anonymous HuggingFace dataset available during review
had been re-uploaded on 18--20 May 2026 with a different
record-level split, which shares 25 of the 174 held-out
records. Until the April manifests and the held-out record
list are published, the segmentation results cannot be
reproduced exactly from public files.
